% 原書369--374頁、Fischler--Susskind mechanism から Lie algebra まで。
\gloss{Fischler--Susskind機構（Fischler--Susskind mechanism）}
{世界面の量子場理論の発散・アノマリーを、弦の摂動論の高次で小さな位相的構造について積分することから生じる発散・アノマリーと相殺すること。量子補正された背景における弦の摂動論の整合性に必要である。}

\gloss{固定点（fixed point）}
{1．幾何では、与えられた対称変換で不変に保たれる点。その変換で空間を同一視すると、境界点または特異点となる。2．量子場理論では、物理が長さの尺度に依存しない、すなわちスケール不変な量子論。通常、このような理論は共形不変でもある。\emph{UV固定点}は量子場理論の短距離の物理を支配する理論、\emph{IR固定点}は長距離の物理を支配する理論である。\emph{自明なIR固定点}は質量のない場を持たない。\emph{非自明なIR固定点}は、ゼロでない相互作用を持つ質量のない場を含む。したがって、IR極限が質量のない自由場理論となる理論は、これらのどちらの用語でも記述されず、\emph{相互作用のないIR固定点}である。}

\gloss{平坦方向（flat direction）}
{スカラー場の空間における、縮退した局所的極小を結ぶ線。この方向に対応する場はモジュラスである。}

\gloss{フロップ（flop）}
{弱結合の弦理論で起こり得る位相の変化。2サイクルが潰れ、その後、別の2サイクルが膨らむ。}

\gloss{分数電荷（fractional charge）}
{電子の電荷の整数倍でない電荷を持つ、閉じ込められていない粒子。ほとんどの$d=4$の弦理論に存在するが、多くの場合、それらはすべて超重粒子である。}

\gloss{分数弦理論（fractional string theory）}
{スピンが$\frac12$の整数倍でない拘束条件を持つ、弦理論の一般化として提案された理論。完全な構成は存在しない。}

\gloss{Fuchs群（Fuchsian group）}
{複素上半平面$H$の$SL(2,\mathbf R)$ Möbius変換の離散部分群$\Gamma$で、$H/\Gamma$が多様体、特にRiemann面となるような付加条件を満たすもの。}

\gloss{汎関数積分（functional integral）}
{本書の用法では、経路積分と同義。}

\gloss{基本領域（fundamental region）}
{$\Gamma$が離散群である商空間$M/\Gamma$に関して、$M$の各点が$F$の内部のちょうど1点、または$F$の境界の1個以上の点と同一視されるような領域$F$。}

\gloss{基本弦、F弦（fundamental string, F-string）}
{D弦やソリトン弦と区別して、その量子化が弱結合の弦理論を定義する、もともとの弦をいう。}

\gloss{融合則（fusion rule）}
{与えられたCFTの任意の2つの一次場の演算子積に、どの共形族が現れるかを指定する規則。}

\gloss{ゲージ固定（gauge-fixing）}
{量子論の冗長なゲージ不変の記述を、各同値類から1つの代表だけを選んだ記述へと縮約すること。}

\gloss{ゲージーノ（gaugino）}
{ゲージボソンと同じ超対称多重項に属するスピン$\frac12$のフェルミオン。}

\gloss{ゲージーノ凝縮（gaugino condensation）}
{ゲージーノ場の積が真空期待値を持つ強結合効果。一般にはカイラル対称性を破るが、直接には超対称性を破らない。しかし、他の場と組み合わさると、しばしば超対称性の破れを誘起する。}

\gloss{世代（generation）}
{クォークとレプトンの1つの組。$SU(3)\times SU(2)\times U(1)$の、カイラルだがアノマリーのない表現の組に属するスピノル場によって記述される。$SU(5)$大統一では$\mathbf5+\overline{\mathbf{10}}$となり、$SO(10)$では$\mathbf{16}$、$E_6$では$\mathbf{27}$に含まれる。\emph{反世代}は共役表現であり、世代と反世代の区別は規約の問題である。}

\gloss{種数（genus）}
{向き付けられた閉Riemann面の取っ手の数$g$。$g=0$は球面、$g=1$はトーラス、以下同様である。}

\gloss{Gepner模型（Gepner model）}
{$N=2$の極小模型CFTに基づく弦模型。}

\gloss{ゴースト（ghosts）}
{Faddeev--Popovゴーストを参照。}

\gloss{Gliozzi--Scherk--Olive射影、GSO射影（Gliozzi--Scherk--Olive projection）}
{フェルミオン場のRおよびNS境界条件について和を取り、世界面フェルミオン数が定まった状態へ射影することによって、モジュラー不変な弦理論を構成する方法。超対称な弦理論では、左向きと右向きのモードに独立なGSO射影がある。両側に同時に作用する対角的な射影は、超対称でない理論を生む。}

\gloss{Goldstino（goldstino）}
{超対称性の自発的破れに伴う、質量のないスピン$\frac12$のGoldstoneフェルミオン。超重力ではグラビティーノと結合し、質量を持つフェルミオンとなる。}

\gloss{Goldstoneボソン（Goldstone boson）}
{自発的対称性の破れの方向の揺らぎに対応する、質量のないスカラー。}

\gloss{大統一（grand unification）}
{$SU(3)\times SU(2)\times U(1)$のゲージ対称性を単純群に統一すること。}

\gloss{大統一スケール（grand unification scale）}
{大統一群が自発的に破れる質量スケール。陽子の安定性と結合の統一から、重力スケールとの差は2、3桁以内でなければならない。}

\gloss{Grassmann変数（Grassmann variable）}
{$\theta_i\theta_j=-\theta_j\theta_i$という関係を満たす代数の元$\theta_i$。フェルミオン場の経路積分表示を与え、また超空間を定義するために使われる。反交換するc数とも呼ばれる。}

\gloss{重力スケール（gravitational scale）}
{無次元の重力結合が1程度になる質量スケール、
$m_{\mathrm{grav}}=\kappa^{-1}=2.4\times10^{18}\,\mathrm{GeV}$。
これはPlanck質量の$(8\pi)^{-1/2}$倍である。}

\gloss{グラビティーノ（gravitino）}
{重力子と同じ超対称多重項に属するスピン$\frac32$のフェルミオン。}

\gloss{Green--Schwarz機構（Green--Schwarz mechanism）}
{Chern--Simons項の$p$形式ポテンシャルの変換則を修正することで、アノマリーを相殺すること。}

\gloss{Green--Schwarz超弦（Green--Schwarz superstring）}
{超対称な弦理論の、超対称性が顕在的な定式化。$\kappa$対称性と呼ばれる、時空のフェルミオン的なゲージ不変性を持つ。簡単な共変的ゲージ固定はない。}

\gloss{$H$モノポール（H-monopole）}
{反対称テンソルゲージ場$B_{\mu n}$の磁荷を持つモノポール。}

\gloss{ヘテロティック5ブレーン（heterotic 5-brane）}
{ヘテロティック弦の質量のない2形式ポテンシャルの磁荷を持つ5ブレーン。ヘテロティック弦のゲージ場のインスタントンの大きさをゼロにした極限として得られる。インスタントン配位は4つの空間次元に局在するので、9空間次元では5ブレーンである。}

\gloss{ヘテロティック弦（heterotic string）}
{左向きと右向きの場に異なる拘束条件代数が作用する弦。現象論的に重要な場合には、$(0,1)$超共形拘束条件代数を持ち、時空超対称性は右向きのモードだけに作用し、ゲージ群は$E_8\times E_8$または$SO(32)$である。}

\gloss{Hagedorn温度（Hagedorn temperature）}
{高励起の弦の状態密度が指数関数的に増大するため、自由弦の熱的分配関数が発散する温度。}

\gloss{隠れたセクター（hidden sector）}
{重力程度の強さの相互作用だけを通じて標準模型と結合する場。隠れたセクターの模型では、超対称性の破れを引き起こす場がこれに含まれる。}

\gloss{最高ウェイト状態（highest weight state）}
{CFTにおいて、すべてのVirasoro下降演算子によって、またはより一般には与えられた代数のすべての下降演算子によって消される状態。}

\gloss{Hodge数（Hodge number）}
{Dolbeaultコホモロジーの非自明な$(p,q)$形式の数。$h^{p,q}$と表す。}

\gloss{ホログラフィック原理（holographic principle）}
{$d$次元の量子重力の状態には、$(d-1)$次元の理論による自然な記述があるという予想。局所場理論からのこの根本的な離脱は、ブラックホールの情報問題に動機付けられ、M理論の定式化の試みで役割を果たしている。}

\gloss{正則（holomorphic）}
{複素変数の理論でいう解析的であること。Minkowski的な接続では左向きになる。反正則な場は共役変数について解析的であり、その接続は右向きである。}

\gloss{正則2次微分（holomorphic quadratic differential）}
{Riemann面上で大域的に定義された正則な$(2,0)$形式。}

\gloss{正則ベクトル場（holomorphic vector field）}
{Riemann面上で大域的に定義された正則な$(-1,0)$形式。}

\gloss{ホロノミー（holonomy）}
{$d$次元多様体上の閉曲線に沿ってベクトルを平行移動すると、元の値を$O(n)$回転した値になって戻ってくる。与えられた多様体について、このようにして得られるすべての回転の集合は$O(n)$の部分群であり、これをホロノミー群という。}

\gloss{ホモロジー（homology）}
{コホモロジーを参照。}

\gloss{超楕円面（hyperelliptic surface）}
{$\mathbf Z_2$対称性を持つRiemann面。分岐切断を持つ球面の2重被覆として表せる。}

\gloss{超Kähler多様体（hyper-Kähler manifold）}
{ホロノミーが$Sp(k)\subset SO(4k)$である$4k$次元多様体。重力が切り離される極限で、$d=6,\ N=1$または$d=4,\ N=2$超対称性のハイパー多重項のモジュライ空間が持つ幾何である。}

\gloss{ハイパー多重項（hypermultiplet）}
{$d=6,\ N=1$または$d=4,\ N=2$超対称性において、ボソン成分が4個の質量のない実スカラーである多重項。}

\gloss{同一視する（identify）}
{2つの点、または他の対象を同値と定義し、それによって商空間を作ること。}

\gloss{赤外発散、IR発散（infrared divergence）}
{時空の長距離から生じる発散。通常、計算すべき量を取り違えたことを意味する。}

\gloss{継承原理（inheritance principle）}
{捩れた理論、すなわちオービフォールド理論において、捩れていない状態の木レベル振幅は、捩れていない理論と同じであるという原理。}

\gloss{挿入（insertion）}
{経路積分の被積分関数から、重み$\exp(iS)$または$\exp(-S)$を除いたもの。}

\gloss{インスタントン（instanton）}
{Euclid経路積分において、作用の大域的最小ではなく局所的最小となる、定数でない配位。通常、このような配位は時空に局在し、位相的に非自明であり、定数配位のまわりの小さな揺らぎからは得られないトンネル効果などを生む場合に重要となる。\emph{時空インスタントン}は時空の有効場理論のインスタントンである。\emph{世界面インスタントン}は世界面の量子場理論のインスタントンであり、時空の非自明な2サイクルに巻き付く世界面に対応する。}

\gloss{交差数（intersection number）}
{一組の曲面が交わる点の数を、交差の向きによって重み付けしたもの。}

\gloss{非関連相互作用（irrelevant interaction）}
{長さの尺度が増すにつれ、その無次元の強さが減少する相互作用。摂動論では、繰り込み不可能な相互作用に等しい。}

\gloss{K3多様体（K3 manifold）}
{4実次元の唯一の非自明なCalabi--Yau多様体。正確には位相的に一意だが、複素構造とKählerモジュライを持つ。ホロノミーは$SU(2)$なので、K3上のコンパクト化により理論の超対称性の半分が破れる。}

\gloss{Kac行列式（Kac determinant）}
{Verma加群の、与えられた$L_0$レベルにある状態の内積行列の行列式。}

\gloss{Kac--Moody代数（Kac--Moody algebra）}
{カレント代数を参照。}

\gloss{Kähler形式（Kähler form）}
{複素多様体上のKähler計量から作る$(1,1)$形式$G_{i\bar\jmath}\,dz^i\,d\bar z^{\bar\jmath}$。}

\gloss{Kähler多様体（Kähler manifold）}
{$n$複素次元でホロノミーが$U(n)$である複素多様体。}

\gloss{Kählerモジュライ（Kähler moduli）}
{Kähler形式をパラメータ化するモジュライ。}

\gloss{Kählerポテンシャル（Kähler potential）}
{Kähler多様体の計量を$G_{i\bar\jmath}=\partial_i\partial_{\bar\jmath}K$によって決めるポテンシャル$K(z,\bar z)$。これは大域的には定義されず、Kähler変換$K(z,\bar z)\to K(z,\bar z)+f(z)+f(z)^*$の任意性を残して定まる。}

\gloss{Kaluza--Kleinゲージ場（Kaluza--Klein gauge field）}
{コンパクト化された理論において、高次元理論の計量に由来するゲージ場。ゲージ群はコンパクト空間の等長変換群である。}

\gloss{Kaluza--Kleinモノポール（Kaluza--Klein monopole）}
{$U(1)$のKaluza--Kleinゲージ対称性の磁荷を持つモノポール。その配位は滑らかなTaub--NUT時空である。3つの空間次元に局在するので、9空間次元では6ブレーンである。}

\gloss{Kaluza--Klein状態（Kaluza--Klein states）}
{コンパクトな空間方向にゼロでない運動量を持つ状態。}

\gloss{Knizhnik--Zamolodchikov方程式、KZ方程式（Knizhnik--Zamolodchikov equation）}
{カレント代数の一次場の期待値を決定する微分方程式。}

\gloss{Landau--Ginzburg模型（Landau--Ginzburg model）}
{ポテンシャルを適切に調整すると非自明なIR固定点を持つスカラー場理論。特に、極小模型CFTのLagrangian表示を与える。}

\gloss{大きな座標変換（large coordinate transformation）}
{恒等変換と連続的につながっていない座標変換。}

\gloss{格子（lattice）}
{$n$次元の線形独立な$n$本の基底ベクトルの、整数係数の線形結合からなる集合$\Gamma$。Euclid計量またはLorentz計量が与えられたとき、\emph{偶格子}とは各点の長さの2乗が偶数である格子である。\emph{双対格子}$\Gamma^*$は、すべての$w\in\Gamma$に対して$v\cdot w\in\mathbf Z$となる点$v$の集合である。\emph{ルート格子}はLie代数のルートの整数係数の線形結合からなる集合である。}

\gloss{レベル（level）}
{1．カレント代数に現れる量子化されたc数項。Schwinger項ともいう。2．自由場理論の振動子励起数の総和。3．共形族において、与えられた状態と最高ウェイト状態の$L_0$固有値の差。}

\gloss{レベル整合（level-matching）}
{$L_0-\widetilde L_0\in\mathbf Z$というモジュラー不変性の条件。}

\gloss{Lie代数（Lie algebra）}
{Jacobi恒等式を満たす反対称な積を持つ代数。\emph{単純Lie代数}は、その補空間と可換な部分代数を持たない。\emph{単純結合Lie代数}では、すべてのルートの長さが等しい。\emph{次数付きLie代数}は奇・偶の元を持ち、奇の元同士の積は対称である。}
